% Skalarprodukt und Winkel – bank/skalarprodukt-und-winkel/ (zusammenbau v0.4, Heft)
\einheitenkopf[t8]{Skalarprodukt und Winkel}
\setcounter{aufgabe}{72}

% Anlauf (ohne Prüfkennung)
\begin{aufgabe}{Das Skalarprodukt als Objekt – Anlauf}
\begin{teile}
\teil $\vec a$, $\vec b$ und $\vec c$ sind Vektoren mit je drei Koordinaten. Was ergibt $\vec a \circ \vec b$? Kreuze an. \\ \kreuz{Vektor} \\ \kreuz{Zahl} \\ \kreuz{nicht definiert}
\teil $\vec a = (2 \mid 1 \mid 3)$, $\vec b = (1 \mid 4 \mid -1)$. Berechne $\vec a \circ \vec b$ und gib an, ob der Winkel zwischen $\vec a$ und $\vec b$ spitz, recht oder stumpf ist. $\vec a \circ \vec b =$ \leerfeld \quad Winkel: \leerfeld
\teil $\vec a$, $\vec b$ und $\vec c$ sind Vektoren mit je drei Koordinaten. Prüfe die Rechenarten von innen nach außen: Was ergibt $(\vec a + \vec b) \circ \vec c$? \\ \kreuz{Vektor} \\ \kreuz{Zahl} \\ \kreuz{nicht definiert}
\end{teile}
\end{aufgabe}

\begin{aufgabe}{Das Skalarprodukt als Objekt – Prüfungsaufgaben}
\begin{teile}
\teil Ein Carsharing-Anbieter verteilt seine $600$ Fahrzeuge auf drei Stadtteile. Der Verteilungsvektor $(x \mid y \mid z)$ gibt die Anzahlen an, es gilt also $x + y + z = 600$. Beurteile beide Aussagen: (1) Es gibt zwei verschiedene Verteilungen, deren Vektoren den Winkel $0^\circ$ einschließen. (2) Es gibt zwei Verteilungen, deren Vektoren den Winkel $90^\circ$ einschließen. \hfill (Abitur 2024 GK)
\teil Ein Quader hat eine quadratische Grundfläche mit der Seitenlänge s und die Höhe h. An einer Ecke der Grundfläche sind $\vec u$ und $\vec v$ die Kantenvektoren der Grundfläche und $\vec w$ der Kantenvektor der Höhe. Begründe ohne Koordinaten, dass $\vec u \circ (\vec u + \vec v + 2\vec w)$ nur von s abhängt. \hfill (Abitur 2021 GK)
\teil Ein Quader hat eine quadratische Grundfläche mit der Seitenlänge s und die Höhe h. An einer Ecke der Grundfläche sind $\vec u$ und $\vec v$ die Kantenvektoren der Grundfläche und $\vec w$ der Kantenvektor der Höhe. $\vec e = \vec u + \vec v$ ist der Vektor der Diagonalen der Grundfläche. Begründe ohne Koordinaten, dass $\vec e \circ (\vec e + 3\vec w)$ nur von s abhängt. \hfill (Abitur 2021 GK)
\end{teile}
\end{aufgabe}

% Anlauf (ohne Prüfkennung)
\begin{aufgabe}{Winkel zwischen Vektoren, Kanten und Geraden – Anlauf}
\begin{teile}
\teil Gesucht ist der Innenwinkel des Dreiecks ABC bei A. Welche zwei Vektoren setzt du an? \\ \kreuz{$\overrightarrow{AB}$ und $\overrightarrow{AC}$} \\ \kreuz{$\overrightarrow{BA}$ und $\overrightarrow{AC}$} \\ \kreuz{$\overrightarrow{AB}$ und $\overrightarrow{BC}$}
\teil Gegeben sind $P(1 \mid 2 \mid 0)$, $Q(3 \mid 3 \mid 2)$ und $R(4 \mid 0 \mid 1)$. Berechne den Winkel zwischen den Kanten PQ und PR. $\varphi \approx$ \leerfeld[$^\circ$]
\teil Gegeben sind $A(2 \mid 1 \mid 0)$, $B(4 \mid 3 \mid 1)$ und $C(5 \mid 1 \mid 3)$. Weise nach, dass das Dreieck ABC rechtwinklig ist, und gib die Ecke mit dem rechten Winkel an.
\end{teile}
\end{aufgabe}

\begin{aufgabe}{Winkel zwischen Vektoren, Kanten und Geraden – Prüfungsaufgaben}
\begin{teile}
\teil Ein Zelt hat die Bodenecken $K(5 \mid 2 \mid 0)$, $L(1 \mid 3 \mid 0)$ und $S(3 \mid 1 \mid 6)$. S ist die Zeltspitze. Berechne die Größe des Winkels zwischen den Kanten KL und KS. \hfill (Abitur 2026 GK) \\ $\varphi \approx$ \leerfeld[$^\circ$]
\teil Ein Würfel mit der Kantenlänge 7 hat eine Ecke im Ursprung, seine Kanten dort liegen auf den positiven Achsen. Auf seinen Kanten liegen $I(7 \mid 0 \mid 1)$, $J(3 \mid 7 \mid 0)$, $K(0 \mid 7 \mid 3)$ und $L(1 \mid 0 \mid 7)$. Das Viereck IJKL ist ein Trapez. Berechne die Größe des Innenwinkels bei I. \hfill (Abitur 2019 GK) \\ $\varphi \approx$ \leerfeld[$^\circ$]
\teil Eine Pyramide ABCDS hat ihre Grundfläche in der xy-Ebene; bekannt sind $C(5 \mid 4 \mid 0)$, $D(1 \mid 8 \mid 0)$ und $S(1 \mid 0 \mid 3)$. Weise nach, dass die Seitenfläche CDS ein rechtwinkliges Dreieck ist. \hfill (Abitur 2024 GK)
\teil Ein Saugroboter fährt vom Punkt $P(1 \mid 5 \mid 0)$ aus geradlinig in Richtung $(4 \mid -1 \mid 0)$ auf eine Wand zu, die am Boden entlang der Strecke von $B(9 \mid 0 \mid 0)$ nach $C(6 \mid 5 \mid 0)$ verläuft; 1 LE entspricht 1 m. Weise nach, dass sich der Roboter der Wand unter einem Winkel von $45^\circ$ nähert. \hfill (Abitur 2021 GK)
\end{teile}
\end{aufgabe}

% Anlauf (ohne Prüfkennung)
\begin{aufgabe}{Neigungswinkel von Ebenen – Anlauf}
\begin{teile}
\teil Der Winkel zwischen zwei Ebenen wird über ihre Normalenvektoren berechnet. Welche Funktion steht in der Formel? \\ \kreuz{Kosinus} \\ \kreuz{Sinus}
\teil Berechne den Winkel zwischen der Ebene $E\colon x + 2y + 2z = 8$ und der xy-Ebene. $\varphi \approx$ \leerfeld[$^\circ$]
\teil Zwei Seitenflächen eines Körpers liegen in den Ebenen $E\colon 2x + y + 2z = 10$ und $F\colon x - 2y + 2z = 3$. Berechne den Schnittwinkel der beiden Ebenen. $\varphi \approx$ \leerfeld[$^\circ$]
\end{teile}
\end{aufgabe}

\begin{aufgabe}{Neigungswinkel von Ebenen – Prüfungsaufgaben}
\begin{teile}
\teil Ein Hang liegt in der Ebene $E\colon 2x + 3y + z = 9$, die xy-Ebene ist die Horizontale. Berechne die Größe des Winkels zwischen Hang und Horizontale. \hfill (Abitur 2025 GK) \\ $\varphi \approx$ \leerfeld[$^\circ$]
\teil Die Pyramide ABCDS hat die Ecken $A(0 \mid 0 \mid 0)$, $B(7 \mid 0 \mid 0)$, $C(7 \mid 4 \mid 0)$, $D(0 \mid 7 \mid 0)$ und die Spitze $S(0 \mid 0 \mid 6)$. Die Seitenfläche BCS liegt in der Ebene $E\colon 6x + 7z = 42$. Berechne den Winkel zwischen E und der Ebene F, in der die Seitenfläche ADS liegt. \hfill (Abitur 2024 GK) \\ $\varphi \approx$ \leerfeld[$^\circ$]
\teil Eine Kletterwand liegt in der Ebene $L\colon 6x + 3y + 2z = 24$, der Untergrund ist die xy-Ebene; 1 LE entspricht 1 m. Berechne die Größe des Winkels zwischen Kletterwand und Untergrund. \hfill (Abitur 2018 GK) \\ $\varphi \approx$ \leerfeld[$^\circ$]
\teil Ein Hausdach liegt in der Ebene $3y - 4z = 0$, die Traufe (untere Dachkante) liegt auf der x-Achse. Von der Traufe aus steigt das Dach über dem Haus (Bereich $y > 0$) an, die Hauswand verläuft von der Traufe aus senkrecht nach unten. Berechne den Neigungswinkel des Dachs gegen die Horizontale und den Winkel im Inneren des Hauses zwischen Dach und Wand. \hfill (Abitur 2024 GK) \\ $\varphi \approx$ \leerfeld[$^\circ$] \quad Innenwinkel \leerfeld[$^\circ$]
\teil Ein Sonnensegel liegt in der Ebene $3y + 8z = 24$, der Untergrund ist die xy-Ebene; 1 LE entspricht 1 m. Damit Regenwasser abläuft, muss die Neigung des Segels mindestens $35\,\%$ betragen. Prüfe, ob diese Bedingung erfüllt ist. \hfill (Abitur 2020 GK)
\teil Eine Kellerklappe ist an der x-Achse drehbar befestigt und steht geöffnet in der Ebene $2y - 3z = 0$; der Boden ist die xy-Ebene. Beim Schließen wird die Klappe gedreht, bis der keilförmige Raum unter ihr kein Volumen mehr hat. Berechne den Drehwinkel und erläutere den Ansatz. \hfill (Abitur 2026 GK)
\teil Eine Zugbrücke ist an der x-Achse drehbar gelagert; hochgezogen liegt sie in der Ebene $3y - z = 0$, die Fahrbahn davor ist die xy-Ebene. Die Brücke wird herabgelassen, bis sie waagerecht aufliegt. Berechne den Drehwinkel und erläutere den Ansatz. \hfill (Abitur 2026 GK)
\end{teile}
\end{aufgabe}

% Anlauf (ohne Prüfkennung)
\begin{aufgabe}{Winkel gegen Ebenen und Bogenmaße – Anlauf}
\begin{teile}
\teil Der Winkel zwischen einer Treppenkante und dem Boden wird über den Richtungsvektor der Kante und einen Normalenvektor des Bodens berechnet. Welche Funktion steht in der Formel? \\ \kreuz{Kosinus} \\ \kreuz{Sinus}
\teil Berechne den Winkel, unter dem die Gerade $g\colon \vec x = (1 \mid 0 \mid 2) + t \cdot (2 \mid 1 \mid 2)$, $t \in \mathbb{R}$, die xy-Ebene schneidet. $\varphi \approx$ \leerfeld[$^\circ$]
\teil Ein Pyramidendach hat die Traufecke $P(3 \mid 3 \mid 2)$ und die Spitze $S(0 \mid 0 \mid 5)$; die Grundfläche des Hauses liegt in der xy-Ebene. Berechne den Neigungswinkel der Dachkante PS gegen die Grundfläche. $\varphi \approx$ \leerfeld[$^\circ$]
\end{teile}
\end{aufgabe}

\begin{aufgabe}{Winkel gegen Ebenen und Bogenmaße – Prüfungsaufgaben}
\begin{teile}
\teil Sonnenlicht fällt in parallelen Geraden mit dem Richtungsvektor $(2 \mid 1 \mid -3)$ auf den ebenen Untergrund, die xy-Ebene. Berechne den Winkel, unter dem das Licht auf den Untergrund trifft. \hfill (Abitur 2023 GK) \\ $\varphi \approx$ \leerfeld[$^\circ$]
\end{teile}
\end{aufgabe}
